\subsection{ELIZA: Conversation by Substitution}
\label{sec:eliza}
\slidesrc{Syntax: RE \& word: 11--15}

ELIZA (Section~\ref{sec:nlp-history}) is the classic proof of how far regular
expressions alone can go. It understands nothing: it applies a \textbf{cascade} of
regex substitutions to the user's sentence and returns the result. A conversation with
an online version (\url{https://www.masswerk.at/elizabot/}) looks like this:

\begin{center}
\begin{minipage}{0.85\linewidth}
\small\ttfamily
ELIZA: How do you do. Please tell me your problem.\\
YOU:\ \ \ I had an unpleasant discussion with a friend of mine\\
ELIZA: Can you elaborate on that ?\\
YOU:\ \ \ She told me that I am stupid\\
ELIZA: Is it because you are stupid that you came to me ?\\
YOU:\ \ \ yes, it is exactly because I am stupid...\\
ELIZA: I see.
\end{minipage}
\end{center}

\emph{``Is it because you are stupid that you came to me?''} is just the user's
\emph{``I am stupid''} with \emph{I am} rewritten to \emph{you are} and pasted into a
template.

\subsubsection{A tiny ELIZA in Python}
\slidesrc{Syntax: RE \& word: 14--15}

Three substitutions applied in cascade with Python's \texttt{re.sub}, which takes a
pattern, a replacement and a string:

\begin{minted}{python}
import re
str = "I feel sad, my heart is broken"
str = re.sub("I", "you", str)        # you feel sad, my heart is broken
str = re.sub("my", "your", str)      # you feel sad, your heart is broken
str = re.sub(r"you (feel|are) (sad|sorry|desperate)(.*)",
             r"Why do you say that you \1 \2\3?", str)
print(str)   # Why do you say that you feel sad, your heart is broken?
\end{minted}

\begin{itemize}[nosep]
  \item Steps 1 and 2 swap the point of view (\emph{I} $\to$ \emph{you},
        \emph{my} $\to$ \emph{your}).
  \item Step 3 uses \textbf{capture groups}: each pair of round brackets in the pattern
        remembers what it matched, and \texttt{\textbackslash 1}, \texttt{\textbackslash 2},
        \texttt{\textbackslash 3} in the replacement paste it back. Here
        \texttt{\textbackslash 1} = \emph{feel}, \texttt{\textbackslash 2} = \emph{sad},
        \texttt{\textbackslash 3} = \emph{, your heart is broken}.
  \item The \texttt{r"..."} prefix makes a \emph{raw string}, so that Python does not
        interpret the backslashes itself.
\end{itemize}

\begin{nbwarning}[order and context matter]
\texttt{re.sub("I", "you", str)} replaces \emph{every} capital \texttt{I}, also inside
words (``\emph{It}'' becomes ``\emph{yout}''), and the order of the cascade matters: if
we swapped \emph{you} $\to$ \emph{I} as well, a later rule could undo an earlier one. Real
ELIZA implementations use word boundaries (\re{\textbackslash bI\textbackslash b}) and
swap pronouns in a single pass.
\end{nbwarning}

\subsubsection{Exercises}
\notinslides

\begin{nbexercise}{1}{prices}
Write a regex matching prices in dollars such as \emph{\$5} and \emph{\$19.99} (an
optional decimal part with exactly two digits). What does it return on ``not \$3.5''?
\end{nbexercise}
\begin{nbsolution}
\re{\textbackslash\$[0-9]+(\textbackslash.[0-9][0-9])?}: an escaped dollar sign (\re{\$}
alone is the end anchor), one or more digits, then an optional group made of a literal
dot and two digits. On ``\$3.5'' the optional group cannot match (only one digit after the
dot), so the regex returns just \emph{\$3}: a false positive of a wrong length.
\emph{Lesson}: a pattern can match only part of what you intended; add
\re{\textbackslash b} or a negative condition after the price if partial matches matter.
\end{nbsolution}

\begin{nbexercise}{2}{reading patterns}
Describe in words what each pattern matches: (a) \re{\textasciicircum[A-Z][a-z]*\textbackslash.\$};
(b) \re{[0-9]+(st|nd|rd|th)}; (c) \re{\textbackslash bcolou?rs?\textbackslash b}.
\end{nbexercise}
\begin{nbsolution}
(a) A whole line made of one capitalised word followed by a full stop (``Stop.'' but not
``Stop now.''). (b) An ordinal written with digits: \emph{1st}, \emph{22nd},
\emph{3rd}, \emph{100th} (and also the wrong \emph{1th}: a regex checks form, not
meaning). (c) The whole words \emph{color}, \emph{colour}, \emph{colors},
\emph{colours}, but not \emph{discolor}.
\end{nbsolution}

\begin{nbsummary}[regular expressions]
\begin{itemize}[nosep]
  \item \re{[abc]} one of; \re{[a-z]} range; \re{[\textasciicircum a]} not (caret first
        in brackets); \re{x|y} one of two strings; \re{(\,)} grouping and capture.
  \item \re{?} 0 or 1; \re{*} 0 or more (Kleene star); \re{+} 1 or more (Kleene plus);
        \re{.} any character; \re{\textbackslash.} a literal dot.
  \item Anchors \re{\textasciicircum} (start), \re{\$} (end), \re{\textbackslash b} (word
        boundary).
  \item Writing a pattern trades off \textbf{false positives} (lower precision) and
        \textbf{false negatives} (lower recall).
  \item ELIZA = a cascade of regex substitutions with capture groups
        (\texttt{\textbackslash 1}, \texttt{\textbackslash 2}, \ldots).
\end{itemize}
\end{nbsummary}
